\chapter{Worked Compilation Example} \label{apx:example}

The following expression combines three operations that the restricted targets must synthesize: a
bitwise XOR, a constant left shift, and a 32-bit constant whose upper bits are set. Compiling it for
rvsc2, rvsc1, and rvsc0 shows how the same source progressively expands as native instructions are
removed from the ISA. The listings below are the real instructions emitted (disassembled with
pseudo-instruction expansion, so \texttt{ret} appears as \texttt{jalr} and \texttt{li} as its
underlying \texttt{lui}).

rvsc1 and rvsc2 both provide \texttt{jalr}, so the operation can be written as an ordinary function
that returns its result:

\begin{lstlisting}[language=C]
int example(int a, int b) {
    return ((a ^ b) << 2) + 0x12345000;
}
\end{lstlisting}

rvsc0 has no \texttt{jalr}, so functions cannot return to a caller. The equivalent rvsc0 program is
a single non-returning function that stores the result to memory instead:

\begin{lstlisting}[language=C]
__attribute__((noreturn))
void example(int a, int b) {
    volatile int r = ((a ^ b) << 2) + 0x12345000;
    __builtin_unreachable();
}
\end{lstlisting}

\textbf{rvsc2} --- full RV32I: every operation is a native instruction.

\begin{lstlisting}
example:
    xor   a0, a0, a1        # a ^ b
    slli  a0, a0, 0x2       # << 2
    lui   a5, 0x12345       # 0x12345000  (low 12 bits are zero)
    add   a0, a0, a5        # + constant
    jalr  zero, 0(ra)       # return
\end{lstlisting}

\textbf{rvsc1} --- no native XOR, shift, or \texttt{jal}/\texttt{ret}. XOR becomes
\texttt{(a|b)-(a\&b)}, the shift becomes repeated self-addition, but \texttt{lui} and \texttt{jalr}
are still native:

\begin{lstlisting}
example:
    and   a5, a0, a1        # \
    or    a0, a0, a1        #  XOR = (a|b) - (a&b)
    sub   a0, a0, a5        # /
    add   a0, a0, a0        # \  << 2 as two self-additions
    add   a0, a0, a0        # /
    lui   a5, 0x12345       # 0x12345000
    add   a0, a0, a5        # + constant
    jalr  zero, 0(ra)       # return
\end{lstlisting}

\textbf{rvsc0} --- additionally lacks \texttt{lui} and \texttt{jalr}. The constant is materialized
with addi+shift (Section~\ref{sec:sc0-lui}) instead of a native \texttt{lui}, and the function stores
its result to the stack rather than returning:

\begin{lstlisting}
example:
    addi  sp, sp, -16
    and   a5, a0, a1        # \
    or    a0, a0, a1        #  XOR = (a|b) - (a&b)
    sub   a0, a0, a5        # /
    add   a0, a0, a0        # \  << 2 as two self-additions
    add   a0, a0, a0        # /
    addi  a5, x0, 72        # hi = 0x12345[19:10]
    [sll  a5, a5, 10]       # a5 = hi << 10
    addi  a5, a5, 837       # a5 = (hi << 10) + lo = 0x12345
    [sll  a5, a5, 12]       # a5 = 0x12345000
    add   a0, a0, a5        # + constant
    sw    a0, 12(sp)        # store result (no return possible)
\end{lstlisting}


\chapter{SLL Synthesis Assembly} \label{apx:sll-asm}

\begin{lstlisting}
# rd = rs1 << rs2
    add   t1, rs2, x0    # t1 = rs2 (save shift count; rd may alias rs2)
    add   rd, rs1, x0    # rd = rs1
    beq   t1, x0, done   # if rs2 == 0, no shift needed
loop:
    add   rd, rd, rd     # rd <<= 1
    addi  t1, t1, -1
    beq   t1, x0, done
    beq   x0, x0, loop
done:
\end{lstlisting}


\chapter{SRL Synthesis Assembly} \label{apx:srl-asm}

\begin{lstlisting}
# rd = rs1 >> rs2 (logical shift right)
    addi  t0, x0, 31
    and   t3, rs2, t0    # t3 = shift & 31
    add   t5, rs1, x0    # t5 = rs1 (save before rd is zeroed)
    addi  rd, x0, 0      # result = 0
    addi  t1, x0, 1      # out_mask = 1
    [sll  t2, t1, t3]    # in_mask = 1 << shift
loop:
    beq   t2, x0, done   # if in_mask == 0, all bits processed
    and   t4, t5, t2     # t4 = rs1 & in_mask
    beq   t4, x0, skip   # if bit is 0, skip
    or    rd, rd, t1     # result |= out_mask
skip:
    add   t1, t1, t1     # out_mask <<= 1
    add   t2, t2, t2     # in_mask <<= 1
    beq   x0, x0, loop
done:
\end{lstlisting}


\chapter{SRA Synthesis Assembly} \label{apx:sra-asm}

\begin{lstlisting}
# rd = rs1 >>_s rs2 (arithmetic shift right)

    # Step 1: check sign bit of rs1 BEFORE srl (avoids rd/rs1 aliasing)
    addi  t0, x0, 1
    [sll  t0, t0, 31]    # t0 = 0x80000000
    and   t6, rs1, t0    # t6 = sign bit (saved in t6; srl uses t0-t5)

    # Step 2: logical right shift
    [srl  rd, rs1, rs2]  # rd = srl(rs1, rs2)

    beq   t6, x0, done   # positive -> no sign extension needed

    # Step 3: re-mask shift; early exit if shift == 0 mod 32
    addi  t0, x0, 31
    and   t3, rs2, t0    # t3 = shift & 31
    beq   t3, x0, done   # shift == 0 mod 32 -> sra(x, 0) = x

    # Step 4: sign_mask = -1 << (32 - shift)
    addi  t2, x0, -1     # t2 = 0xFFFFFFFF
    addi  t4, x0, 32
    sub   t4, t4, t3     # t4 = 32 - (shift & 31)
    [sll  t2, t2, t4]    # sign_mask = -1 << (32 - shift)
    or    rd, rd, t2
done:
\end{lstlisting}


\chapter{SLT Synthesis Assembly} \label{apx:slt-asm}

\begin{lstlisting}
sub   t0, rs1, rs2   # diff = rs1 - rs2
[xor  t1, rs1, rs2]  # t1 = rs1 ^ rs2  (derived)
[xor  t2, rs1, t0]   # t2 = rs1 ^ diff  (derived)
and   t1, t1, t2     # overflow = (rs1^rs2) & (rs1^diff)
[xor  t0, t0, t1]    # corrected = diff ^ overflow  (derived)
[srl  rd,  t0, 31]   # rd = corrected >> 31  (derived)
\end{lstlisting}


\chapter{SLTU Synthesis Assembly} \label{apx:sltu-asm}

\begin{lstlisting}
sub   t0, rs1, rs2   # diff = rs1 - rs2
[not  t1, rs1]       # t1 = ~rs1  (derived)
and   t2, t1, rs2    # t2 = ~rs1 & rs2  (borrow generated)
[xor  t3, rs1, rs2]  # t3 = rs1 ^ rs2  (derived)
[not  t3, t3]        # t3 = ~(rs1 ^ rs2)  (derived)
and   t3, t3, t0     # t3 = ~(rs1^rs2) & diff  (borrow propagated)
or    t2, t2, t3     # borrow = generated | propagated
[srl  rd,  t2, 31]   # rd = borrow >> 31  (derived)
\end{lstlisting}


\chapter{LB / LBU / LH / LHU Synthesis Assembly} \label{apx:lb-lh-asm}

\begin{lstlisting}
# lb rd, 0(rs1)  (address in rs1, byte_pos unknown at compile-time)
addi  t0, x0, -4
and   t0, rs1, t0         # t0 = rs1 & -4  (word-aligned)
lw    t1, 0(t0)           # t1 = word containing the byte
addi  t0, x0, 3
and   t0, rs1, t0         # t0 = rs1 & 3  (byte position: 0-3)
[sll  t0, t0, 3]          # t0 = byte_pos * 8  (bit offset)
[srl  t1, t1, t0]         # t1 >>= bit_off  (byte in bits [7:0])
[sll  t1, t1, 24]         # t1 <<= 24  (byte in bits [31:24])
[sra  rd,  t1, 24]        # rd >>= 24  (sign-extend -> lb)
                          # use [srl] in the last step for lbu (zero-extend)

# lh rd, 0(rs1)  -- identical but MASK = 2, BITS = 16
addi  t0, x0, -4
and   t0, rs1, t0         # word-aligned
lw    t1, 0(t0)
addi  t0, x0, 2
and   t0, rs1, t0         # t0 = rs1 & 2  (0 or 2: which halfword)
[sll  t0, t0, 3]          # t0 = half_pos * 8  (0 or 16)
[srl  t1, t1, t0]         # halfword in bits [15:0]
[sll  t1, t1, 16]
[sra  rd,  t1, 16]        # sign-extend -> lh  (srl for lhu)
\end{lstlisting}


\chapter{SB Synthesis Assembly} \label{apx:sb-asm}

\begin{lstlisting}
# sb rs2, 0(rs1)   -- address in rs1, byte position unknown at compile-time
addi  t0, x0, -4
and   t3, rs1, t0           # t3 = rs1 & -4  (word-aligned)
addi  t0, x0, 3
and   t0, rs1, t0           # t0 = rs1 & 3  (byte_pos: 0-3)
add   t0, t0, t0            # \
add   t0, t0, t0            #  t0 = byte_pos * 8  (constant-3 sll; no loop needed)
add   t0, t0, t0            # /
lw    t1, 0(t3)             # t1 = old word
addi  t2, x0, 255           # t2 = 0xFF
[sll  t2, t2, t0]           # t2 = 0xFF << shift  (mask)
[not  t2, t2]               # t2 = ~mask  -- 2 insns
and   t1, t1, t2            # t1 = old_word & ~mask  (clear target byte)
addi  t2, x0, 255
and   t2, rs2, t2           # t2 = rs2 & 0xFF  (isolate input byte)
[sll  t2, t2, t0]           # t2 = byte_value << shift
or    t1, t1, t2            # t1 = word with byte inserted
sw    t1, 0(t3)             # write back
\end{lstlisting}


\chapter{SH Synthesis Assembly} \label{apx:sh-asm}

\begin{lstlisting}
# sh rs2, 0(rs1)
addi  t0, x0, -4
and   t3, rs1, t0           # t3 = rs1 & -4  (word-aligned)
addi  t0, x0, 2
and   t0, rs1, t0           # t0 = rs1 & 2  (hw_pos: 0 or 2)
add   t0, t0, t0            # \
add   t0, t0, t0            #  t0 = hw_pos * 8  (0 or 16; constant-3 sll)
add   t0, t0, t0            # /
lw    t1, 0(t3)             # t1 = old word
lui   t2, 0x10              # \
addi  t2, t2, -1            #  t2 = 0xFFFF
[sll  t2, t2, t0]           # t2 = 0xFFFF << shift  (mask)
[not  t2, t2]               # t2 = ~mask  -- 2 insns
and   t1, t1, t2            # t1 = old_word & ~mask
lui   t2, 0x10
addi  t2, t2, -1            # t2 = 0xFFFF
and   t2, rs2, t2           # t2 = rs2 & 0xFFFF
[sll  t2, t2, t0]           # t2 = hw_value << shift
or    t1, t1, t2            # t1 = word with halfword inserted
sw    t1, 0(t3)             # write back
\end{lstlisting}
